\section{Object storage, documents and fine-grained access}\label{sec:content}
Part~3 defines the storage product delivered in R1: a native object API, a canonical namespace and version catalog, durable publication, and the fine-grained access controls shared by objects, documents and text units. HTTP clients and applications can use this product with no MCP server, tool executor or model provider deployed. R1-B delivers the object service and R1-C adds document and text projections over the same objects. Every payload is stored encrypted under its bucket's mode (\rr{ENC-001}): in server mode, the default, the storage role encrypts uploads and decrypts for authorized readers under data keys released by the \code{keys} role; in client mode, clients encrypt before upload and decrypt after authorized download, and the server stores ciphertext and opaque recipient envelopes only. R1-C gives independently protected units separate keys in both modes. Part~4 defines both modes and their key ownership. Backend encryption at rest remains an additional layer. MCP is a later service client of this foundation, not its storage or authorization authority.

\subsection{Storage application service and native API}
\begin{lstlisting}
// Logical application interface; all callers retain verified identity context.
type ContentService interface {
    Read(ctx context.Context, req ReadContentRequest) (ContentPackage, error) // plaintext (server mode) or encrypted package (client mode)
    Stat(ctx context.Context, req StatContentRequest) (ContentMetadata, error)
    List(ctx context.Context, req ListContentRequest) (ContentPage, error)
    Write(ctx context.Context, req WriteContentRequest) (PublishedVersion, error)
    Delete(ctx context.Context, req DeleteContentRequest) (DeleteOutcome, error)
    Export(ctx context.Context, req ExportContentRequest) (ExportHandle, error)
}
\end{lstlisting}
\req{CNT-001}{R1-B}{T-06}{\code{internal/content} is the single application layer for content reads, stat and discovery, listing, writes, deletion, authorized envelope delivery and server-mode key release. The native HTTP API delegates to it in R1, the MCP service calls its authenticated API in R2, and S3 handlers delegate to it in R3. It resolves canonical aliases and versions and performs the business checks and mandatory guards. A client-mode read returns an encrypted package: exact ciphertext, its bound signed descriptor or authorized manifest projection, and an existing envelope for the authorized recipient under Table~\ref{tab:keyops}; the service never receives client-mode payload keys or decrypts client-mode payloads. A server-mode read returns plaintext over the authenticated connection, decrypted by the storage role with a data key released under \rr{SSE-003}.}
\req{CNT-002}{R1-B}{T-06}{Protocol adapters translate authentication evidence, selectors, operation parameters and response framing. They \mustnot\ perform a permissive content check and then bypass the shared service.}
\req{CNT-003}{R1-C}{T-06}{For a document registered with protected units, an original-file path authorizes its complete disclosure set. Parent permission or a permitted section cannot authorize a forbidden section through the encrypted original, its envelope or a server-mode original read. The original is a separate encrypted object with a separate key; delivery requires the authoritative full-content binding and every required decision. Selected unit keys never decrypt the original. Authorized-subset projection is explicit and does not silently change an ordinary object download.}
\req{CNT-004}{R1-B}{T-09}{HEAD and error responses \mustnot\ reveal existence or metadata to callers that lack stat or discovery authority.}
\req{CNT-005}{R1-B}{T-09}{A listing authorizes bucket listing, scans committed catalog candidates in stable key order, batch-checks per-object discovery, and accumulates visible entries until the page size or a bounded scan or time budget is reached. A short or empty page may carry a continuation token and the cursor must make progress. Metadata is filtered as well as names. A continuation token binds tenant, principal, identity revision, query, prefix and delimiter, policy epoch, recovery incarnation, expiry and last scanned position; its contents are encrypted and authenticated or held in opaque server state, and replay under another identity or a changed query fails. \code{CommonPrefixes} \mustnot\ expose hidden children, and delimiter support stays disabled until safe prefix discovery passes tests.}
\req{CNT-006}{R1-B}{T-09}{Listings and discovery are read-only from the caller's perspective. They may report a bounded integrity hint but \mustnot\ execute expiration, deletion, declassification or replication as a hidden side effect; background maintenance owns those effects under its own authority and job record.}
\reqn{CNT-007}{R1-B}{T-09, T-22}{Candidate selection for listing and search may use a sound superset prefilter (for example tenant and project predicates evaluated in SQL). Every candidate that survives still receives an exact per-resource decision, and a prefilter is never the only check.}
\req{CNT-008}{R1-B}{T-13}{A read admits its exact plan under \rr{DAT-007}, with a session check only if its profile requires one. Native reads and baseline MCP need no agent session. An admitted stream compares all bound authority and transitive lineage dependencies at least every $T_{check}$; the initial target is 5 s, qualified under T-22. Failed or overdue checks (no completion within $2T_{check}$) abort the connection rather than report normal completion. A separate deadline prevents new emission at or after \code{valid_until}, including a temporal policy cutoff; it does not wait for the next periodic revision check. Continuing requires fresh admission. Cancellation ends the lease. Already-emitted plaintext, ciphertext and envelopes and retained client keys remain outside revocation.}

The native HTTP API is a versioned typed contract, not a set of MCP tools. It accepts bounded plaintext for a server-mode bucket, which the storage role encrypts, and the qualified encrypted representation of arbitrary bounded object bytes for a client-mode bucket. In client mode the client SDK encrypts binary files without parsing them and client-side document ingestion splits text into independently encrypted units; in server mode the storage role splits uploaded text under the same profile (\rr{DOC-004}). An ordinary binary object is authorized as a whole object. Once it is registered as a document with independently protected units, its original-file route must satisfy \rr{CNT-003}. Table~\ref{tab:nativeops} is the minimum supported operation set. It names domain operations rather than prescribing URL spelling; the OpenAPI contract freezes the request, response and error mapping before qualification.
\begin{longtable}{@{}P{0.20\textwidth}P{0.43\textwidth}P{0.25\textwidth}@{}}
\caption{Native storage operations, available before MCP}\label{tab:nativeops}\\
\toprule Operation & Shared authorization and publication boundary & Availability\\\midrule\endfirsthead
\toprule Operation & Shared authorization and publication boundary & Availability\\\midrule\endhead
Bucket administration & Authorized control API creates the tenant-bound namespace and protected defaults; clients cannot set their own authority metadata & R1-B\\
Create / write object & Bucket create-object or existing-object new-version; bounded plaintext encrypted by the storage role (server mode) or client-encrypted upload with signed descriptor (client mode); fresh commit decision and conditional publication & R1-B\\
Read / stat object & Resolve the exact committed version, enforce read or stat and return plaintext (server mode) or authorized ciphertext and recipient envelope (client mode), or permitted metadata; historical access is explicit & R1-B\\
List objects / versions & Bucket listing plus per-entry discovery, protected metadata projection and an opaque continuation; historical versions need their own permitted visibility & R1-B\\
Delete key / version & Authorized logical tombstone or explicit version deletion; retention and references checked before physical retirement & R1-B\\
Share grants & \code{share} on the resource records or removes a grant; a grant takes effect only where a rule honors it (\rr{SHR-001}, \rr{SHR-002}) & R1-B\\
Read document / text & Parent and selected-unit checks over the same stored objects; the original file and its separate key require complete-content permission & R1-C\\
Discover / search & Server filters permitted catalog metadata; a client searches plaintext locally after an authorized read (and, in client mode, envelope delivery and decryption) & R1-C\\
\bottomrule
\end{longtable}
\reqn{CNT-009}{R1-B}{T-06, T-08, T-09}{The native storage API implements the R1-B operations in Table~\ref{tab:nativeops} through \code{ContentService} and the authorized control API. Its versioned contract fixes object-size and request bounds, stable object and version selectors, required actions, conditional-write behavior, idempotency, cancellation, and redacted errors. Each read and mutation binds the exact tenant, canonical resource and immutable version or expected generation. Object CRUD and listing remain usable when every MCP, model and tool-execution service is absent; clients receive no backend credentials or direct backend download URL that bypasses these checks.}

\subsubsection{Native version and visibility semantics}
These mappings specialize \rr{RES-003}, \rr{RES-006} and \rr{CNT-009}; they are shared service behavior, not additional adapter policies.
\begin{longtable}{@{}P{0.23\textwidth}P{0.65\textwidth}@{}}
\toprule Request & Exact operation meaning\\\midrule\endhead
Latest read / stat & Resolve the current pointer once to a committed immutable version; check \code{read} or \code{stat} on that object/version. An encrypted-package read also needs every recipient-envelope check in Table~\ref{tab:keyops}. A delete marker means no current content. A later overwrite never makes an admitted stream switch bytes.\\
Historical read / stat & Naming a non-current version additionally requires \code{read-history} for its bytes, or \code{discover-history} for its metadata, together with \code{read} or \code{stat}. The plan binds that exact version and its current protected attributes. Envelope delivery also needs its recipient checks in Table~\ref{tab:keyops}; knowing a version identifier grants nothing.\\
Version listing & Require bucket \code{list}; disclose each entry only with its \code{discover} check and, for a non-current version, \code{discover-history}. Omit unauthorized version identifiers, checksums and timestamps, including delete markers; return no total version count.\\
Create / overwrite & For a new key require bucket \code{create-object}; for an existing object require \code{new-version}. Client-selected owner, tenant and label fields cannot set protected attributes. An overwrite creates a version; it never modifies published bytes.\\
Key / version deletion & Key deletion requires \code{delete} and creates a current delete marker. Permanent deletion requires \code{delete-version} on the exact version plus retention and reference guards. Deleting the current version leaves a delete marker; it never silently restores an older version.\\
\bottomrule
\end{longtable}

Preconditions do not disclose hidden state. Authenticate and check the required operation before returning whether an ETag matched, an object exists, or a version has already been deleted. The native profile uses one generic not-found response for an absent or undiscoverable target; a caller with discovery but insufficient operation authority may receive a generic denial. Both omit protected identifiers and distinguish dependency failure from success. \code{If-Match} compares an opaque ETag for the exact logical version; \code{If-None-Match: *} means no current live version at commit, so a tombstoned key may receive a new version after authorization. Preconditions apply to the catalog transaction, never to an independently observed backend object. The ETag identifies the immutable encrypted representation, never a plaintext hash. No ETag or ciphertext digest is returned without the corresponding metadata right.

The baseline transfers complete encrypted object or unit streams under \rr{CRY-007}; range requests are rejected until a separately qualified ciphertext-range profile exists. Such a future profile must authenticate all required format headers and segments on the client and preserve whole-object or whole-unit authorization. A range never creates a paragraph-level right or bypasses document-unit checks. A request for explicit document units either authorizes the entire requested set or denies; a separately named subset operation returns only permitted units and metadata under \rr{DOC-007}. Responses advertise the selected version, representation length and completion semantics only after authorization. If ciphertext transport integrity, lease validation or a backend read fails after headers or bytes have been sent, the server aborts the stream rather than returning a normally completed shorter object. Clients also reject failed AEAD authentication or truncation; the server cannot verify payload AEAD tags. An incomplete transfer is a failure and is retried against the same version under fresh authority.

\subsection{Catalog records and ownership}
The catalog stores logical identity, protected policy attributes and publication state; the backend stores immutable ciphertext. Required policy labels and minimal routing metadata are visible to the authorized server roles. Substantive titles, source text and other sensitive descriptive content remain inside encrypted client payloads; opaque aliases avoid leaking them through object keys. Table~\ref{tab:storagecatalog} specializes \rr{DAT-001}--\rr{DAT-004}. Every primary, unique and foreign-key reference below includes the tenant, even when an identifier is globally generated. Queries also carry the tenant predicate; identifier unpredictability is not isolation.
\begin{longtable}{@{}P{0.19\textwidth}P{0.43\textwidth}P{0.26\textwidth}@{}}
\caption{Storage catalog records and uniqueness}\label{tab:storagecatalog}\\
\toprule Record & Stored fields and invariant & Writer\\\midrule\endfirsthead
\toprule Record & Stored fields and invariant & Writer\\\midrule\endhead
Bucket / alias & Bucket UUID, namespace name, protected defaults and revision. Unique tenant/name. Object alias is unique tenant/bucket/exact-key and points to a stable object UUID. Names never become backend paths. & Authorized namespace/control operation\\
Object head & Stable object UUID, current version or delete marker, generation. At most one current pointer; changing aliases preserves permission identity. Referenced version belongs to this tenant and object. & Publication or logical deletion transaction\\
Immutable version / blob & Unique tenant/object/version; blob identity, ciphertext length and digest, client-crypto profile, signed descriptor, durable receipt and lifecycle state. Payload identity and bytes never change; a state change cannot retarget a version to other bytes. & Upload finalization and publication; retirement changes state only\\
Upload intent & Unique tenant/subject/operation-class/idempotency-key; actor/delegation binding, canonical argument and ciphertext digests, expected generation/version, assigned logical-version, crypto-version and blob IDs, encrypted-attempt ID, state and stored result. & Storage operation owner; reconciler resumes by this identity\\
Protected attributes & Object or unit owner, project, label and independent metadata/access revisions. Content generation and access revisions are not interchangeable. & Explicit authorized metadata/label operation\\
Manifest / references & Unique tenant/document/version and unit membership; uploader-signed manifest digest, exact unit versions and ciphertext hashes. Full structure is restricted metadata; subset responses use signed per-unit descriptors. Foreign keys prevent cross-tenant references; live reference edges protect units. & Document publication transaction in R1-C\\
\bottomrule
\end{longtable}

\textbf{Transaction ownership.} One catalog transaction function owns each publication, metadata change or logical deletion and its audit/outbox append. Callers pass a prepared immutable plan and expected revisions; they cannot supply arbitrary SQL or set lifecycle fields directly. Network uploads, provider calls and policy evaluation run outside database write transactions. The final transaction rechecks the prepared decision and state preconditions under \rr{DAT-007} before changing visibility. The catalog and backend do not use a distributed transaction: a finalized but unpublished blob is recoverable garbage, while a published version must already have its qualified durable receipt.

\textbf{Concurrency discipline.} The publication design records one consistent acquisition order covering authority, namespace, versions, intents and references. Document publication, garbage collection and reference deletion obey the same protocol. No provider or policy RPC runs while database locks are held. A failed conditional transition is a conflict, never success; bounded retries preserve operation identity and refresh stale authority.

\textbf{Read sequence.} Resolve the alias and requested version from authoritative catalog state; load trusted object, identity and delegation attributes; obtain the exact signed checks; then atomically revalidate those bindings and record admission under \rr{DAT-007}. Open only the recorded immutable ciphertext through the backend adapter. Independently authorize \code{GetRecipientEnvelope} under Table~\ref{tab:keyops} and return only the matching opaque envelope and signed descriptor. The recipient client verifies signatures, exact object/version/unit/recipient bindings and the ciphertext format, unwraps with its locally held private key, then authenticates and decrypts the payload locally. Hold the bounded stream lease of \rr{CNT-008}. Backend errors never trigger a fallback to a different version, a raw download URL or a less restrictive policy. Concurrent retirement may make a read fail or abort; it cannot substitute other content.

\textbf{Write and retry sequence.} Reserve an upload intent and allocate object, logical-version, crypto-version and blob identities before constructing the client encryption context. The client creates the bound ciphertext, signed descriptor and separate recipient envelopes and supplies the bounded ciphertext length and digest. Server validation checks signer authority, identifier bindings, format declarations, length and digest; it cannot validate plaintext, AEAD tags or label accuracy. Initial envelopes require the same complete sharing and target-recipient check sets as later \code{PublishRecipientEnvelope}, not merely create-object permission (\rr{KEY-012}, \rr{RES-009}). Finalize the backend receipt before the short publication transaction of \rr{PUB-003}. A retry of the same encrypted attempt resends identical ciphertext, descriptors and envelope bytes with identical IDs. Restarting encryption requires a new upload intent, encrypted attempt and idempotency identity, with a fresh logical version, keyset, crypto version and blob identity under the same stable object. Old immutable ciphertext is never overwritten, and fresh signatures bind all new identities. A completed exact retry returns the original result only after authenticating the owner and checking current permission to observe it; it never republishes. An interrupted finalize is reconciled by stored attempt/blob identity, and an uncertain catalog commit follows \rr{AVL-004}, not mere visibility of an operation row.

\textbf{Delete and reclamation sequence.} Authorize the exact key or version, lock the same rows, recheck retention, current references and revisions, then publish the logical withdrawal and audit record. A durable retirement record schedules physical removal only after that transition. The worker rechecks current references and its retirement authority immediately before deletion; retries are idempotent by tenant/blob identity. It never infers abandonment from file age or lists the backend as a source of user-visible objects. Restore and orphan reconciliation use catalog identity and receipts under \rr{PUB-006}; unmatched data remains quarantined until reconciled.

\subsection{Publication protocol}\label{sec:publication}
PostgreSQL controls logical namespace and visibility. Public bucket names and object keys are never used as filesystem paths.
\begin{figure}[ht]
\centering
\begin{tikzpicture}[st/.style={draw=ink!70,rounded corners=2pt,fill=pale,align=center,minimum height=8mm,minimum width=22mm,font=\small},arr/.style={-{Stealth[length=2mm]},ink}]
\node[st] (res) {RESERVED};
\node[st,right=10mm of res] (wr) {WRITING};
\node[st,right=10mm of wr] (stg) {STAGED};
\node[st,right=18mm of stg] (live) {LIVE};
\node[st,below=12mm of live] (withdrawn) {WITHDRAWN};
\node[st,left=18mm of withdrawn] (retired) {RETIRED};
\node[st,below=12mm of wr] (ab) {ABORTED /\\QUARANTINED};
\draw[arr] (res) -- (wr);
\draw[arr] (wr) -- (stg);
\draw[arr] (stg) -- node[above,font=\scriptsize]{publish} (live);
\draw[arr] (live) -- node[right,font=\scriptsize]{delete version} (withdrawn);
\draw[arr] (withdrawn) -- node[above,font=\scriptsize]{reclaim} (retired);
\draw[arr] (res.south) |- (ab.west);
\draw[arr] (wr) -- (ab);
\draw[arr] (stg.south) -- (ab.north east);
\end{tikzpicture}
\caption[Version lifecycle]{Key deletion changes the object head to a separate delete marker; it does not change an older version's state. Version withdrawal and physical retirement require \rr{PUB-005}. Staged records may be durable without being publicly readable.}\label{fig:pubstates}
\end{figure}
\req{PUB-001}{R1-B}{T-08}{PostgreSQL controls logical namespace and visibility. Public bucket names and object keys \mustnot\ be used directly as filesystem paths.}
\req{PUB-002}{R1-B}{T-08, T-11}{Staging proceeds as follows. (1) \code{BeginEncryptedUpload} creates an intent with tenant, actor, subject, delegation, assigned object, logical-version, crypto-version and blob IDs, encrypted-attempt ID, expected generation/current version, metadata revision and idempotency key. (2) Authorize create or new-version and assign protected metadata, including the bucket's encryption mode. (3) In client mode the client constructs the exact bound ciphertext, signed descriptor and one independent envelope per object or unit and recipient; in server mode the storage role encrypts the uploaded plaintext under \rr{SSE-001} and \rr{SSE-005} with a data key obtained under \rr{SSE-003}, and there are no client descriptors or envelopes. (4) \code{FinalizeEncryptedUpload} validates exact identifiers, format, ciphertext digest/length and, in client mode, signatures and envelope references, obtains the backend durable receipt, and requires every initial envelope to satisfy the same caller-share/create and target-recipient read/receive-key checks as \code{PublishRecipientEnvelope} under \rr{KEY-012} and \rr{RES-009}. Target eligibility uses the trusted intended committed state within this upload plan, never a staged-data read grant. The version and its envelopes remain unpublished until the complete fresh commit succeeds. No client-mode upload route accepts a payload decryption key.}
\req{PUB-003}{R1-B}{T-08, T-11}{Commit obtains a fresh decision for object publication and every initial envelope. Under \rr{DAT-007}, one transaction checks expected generation/current version, metadata, caller and target-principal revisions, recipient-key state, delegation, policy epoch, recovery incarnation, decision expiry and conditional-write preconditions. A missing, stale or misbound check aborts. It atomically publishes the version and approved envelopes, updates the head, consumes the intent and appends audit/outbox records. The same boundary publishes newly staged document units under \rr{DOC-002}. Initial envelopes retain all \rr{KEY-012} checks; success requires the durability barrier.}
\req{PUB-004}{R1-B}{T-08}{Only a committed LIVE version is eligible for ordinary reads. A durable STAGED catalog record, finalized backend blob or failed commit grants no visibility; pending versions and envelopes remain inaccessible through every route.}
\req{PUB-005}{R1-B}{T-08}{Key deletion publishes a separate current-head delete marker and advances the head generation; historical LIVE versions remain readable only with historical rights. Explicit version deletion checks version-delete permission, retention and references, then durably changes that version to WITHDRAWN and schedules reclamation. Deleting the current version also leaves a head marker, never restores an older version. Physical removal follows the durable withdrawal and fresh retirement guards; only confirmed removal marks RETIRED.}
\req{PUB-006}{R1-B}{T-08}{Garbage collection uses catalog references, intents, withdrawals and a grace period; file age alone is insufficient. Reconciliation is idempotent across blob finalization and catalog commit. Expired RESERVED, WRITING or STAGED attempts can be aborted under expected-state checks. Restore reconciles catalog and blobs or quarantines mismatches, and reopens access only under \rr{AVL-008}.}
\req{PUB-007}{R1-B}{T-08}{\code{If-Match} and \code{If-None-Match} apply to the committed logical version and its documented ETag, not to a staging file. Conflicting commits serialize with row locks or compare-and-swap and a fixed lock order, and two concurrent overwrites \mustnot\ both publish against the same obsolete expected generation. An ETag is an opaque identifier and is not promised to be an MD5; the server stores ciphertext checksums only, in either mode. A plaintext digest, if needed by a client, stays inside its encrypted metadata and is verified locally.}
\req{PUB-008}{R1-B}{T-08}{An upload follows RESERVED, WRITING, STAGED and LIVE; any pre-LIVE state may end in ABORTED or QUARANTINED. Explicit version deletion follows LIVE to WITHDRAWN to RETIRED. Delete markers belong to the object-head history, not this version lifecycle. Transitions use expected-state checks, and no state change retargets immutable bytes (Figure~\ref{fig:pubstates}).}
\req{PUB-009}{R1-B}{REV}{The first release does not claim regulatory WORM or Object Lock compliance.}

\subsection{Object backend contract}
\begin{lstlisting}
// Logical contract. All requests carry verified workload bindings.
type ObjectBackend interface {
    Begin(ctx context.Context, req BeginRequest) (StagedWriter, error)
    Finalize(ctx context.Context, req FinalizeRequest) (DurableReceipt, error)
    Open(ctx context.Context, req ReadRequest) (BlobReader, error)
    Remove(ctx context.Context, req RetireRequest) (RemovalReceipt, error)
}
\end{lstlisting}
\req{BKD-001}{R1-B}{T-08, T-15}{\code{Begin} binds operation, tenant, blob and placement. \code{Finalize} checks exact length and digest and returns a durable receipt only after the backend's qualified durable completion condition: for the R1 single-node store, the required file and directory flushes after an exclusively created temporary file is fully written and verified; for a later backend, its own qualified acknowledgement bound to the immutable object identity and verified content. A successful network send alone is not a durable receipt. Repeating the same identity and content is idempotent, and the same blob ID with different bytes rejects. A timeout is an unknown result reconciled by operation identity, not an assumed failed write. \code{Open} exposes exact immutable bytes, and \code{Remove} accepts only a current catalog-authorized retirement plan. The caller owns stream close, cancellation and logical publication.}
\reqn{BKD-002}{R1-B}{T-08}{R1 runs on our own single-node blob store on local disk behind this contract (decision D-02). It writes each blob to an exclusively created temporary file, checks length and digest, flushes the file and its directory, and only then returns the durable receipt; the catalog, not the file system, decides visibility. It receives only ciphertext, from the client in client mode or from the storage role's encryption in server mode, and no other role has direct file access. R1 makes no zero-loss claim for blob content (N-08): losing the node loses blobs written since the last backup. A later backend (the replicated store of Part~5, or another store qualified under \rr{BKD-004}) replaces it without changing client policy semantics. Volume encryption supplements payload encryption and gives no platform role a payload decryption key.}
\req{BKD-003}{R1-B}{T-08}{The local backend uses service-owned directories, restrictive permissions, descriptor-relative safe file operations and symlink and traversal protection. Finalized blobs are immutable, and file and required directory metadata are flushed before a catalog transaction publishes a LIVE reference. Storage clients have no direct filesystem or backend credentials.}
\req{BKD-004}{R1-B}{T-08}{A different backend implements the same interface and needs its own consistency and durability tests. It does not duplicate user-facing policy combination in a second storage engine.}

\reqn{BKD-005}{R1-B}{T-08, T-18}{Payload encryption is mandatory from R1-B in both modes (\rr{ENC-001}). A server-mode upload is encrypted by the storage role before any byte reaches a staged blob, temporary file or spill (\rr{SSE-005}). A client-mode upload declares the approved encrypted format, supplies its signed binding and is produced by a qualified encrypting client. Payloads remain ciphertext in unpublished uploads, temporary files, spill, finalized blobs, replicas and backups, and no route stores payload plaintext; no server path decrypts a client-mode payload. Backend and volume encryption also protect persistent server records, opaque envelopes, wrapped keys and ciphertext, with an approved configuration, verification method and documented drift-detection bound. For client mode, a declaration or signature proves origin and binding, not that arbitrary maliciously supplied bytes were encrypted correctly; the cryptographic guarantee assumes the qualified client follows the format, and the server cannot prove payload AEAD validity without a key and does not claim to do so. Unsupported or missing profiles reject publication.}

\subsection{Documents and text units in the storage core}
Native document and text-unit endpoints are part of R1-C in both modes. They call the same resolver and delivery path later used by MCP; a document is a catalog projection over independently protected encrypted units, not an MCP-owned resource.
\req{DOC-001}{R1-C}{T-06, T-08}{A document version has a canonical manifest identifying its exact unit membership, unit/object versions, ciphertext digests and format, and declared extraction and normalization versions. Its digest is computed over the canonical core of unit entries before per-unit descriptors or signatures are added, avoiding self-reference. In client mode the manifest is uploader-signed, and each deliverable unit has an uploader-signed descriptor binding tenant, document/version, the opaque manifest digest, unit/version, ciphertext digest and crypto profile; the server verifies signatures and cross-references at publication, and a subset client verifies its descriptors without receiving the hidden full manifest. In server mode the storage role builds the manifest and unit bindings from the uploaded text and records them in the catalog with the authenticated uploader; there is no client signature.}
\req{DOC-002}{R1-C}{T-08}{New document units remain STAGED, with durable receipts and nonpublic catalog records, until the document transaction atomically publishes their LIVE versions, permitted envelopes, manifest and protective references. Already-published units may be referenced after current checks. In client mode a reused unit needs new document-specific descriptors and recipient envelopes prepared by an authorized client holding the unit key, and the server cannot create missing envelopes; in server mode reuse needs no client key, and the storage role binds the reused unit to the new manifest. Reuse preserves the immutable unit ciphertext. Every referenced unit must be LIVE at the transaction outcome, and incomplete ingestion is never exposed through raw, envelope or plaintext APIs. The parent and all affected units use the same authority fence and publication preconditions; retirement honors their references.}
\req{DOC-003}{R1-C}{T-06}{The qualified client in client mode, or the storage role in server mode, splits content at authorization boundaries into stable sections, paragraphs or independently protected entries and encrypts each unit with a separate payload key. Re-splitting creates a new document version. Protected access labels and revisions remain separate mutable server records; label changes alone do not change payload bytes, confer a missing recipient envelope or release a key.}
\req{DOC-004}{R1-C}{T-09}{The initial ingestion profile accepts bounded UTF-8 text and Markdown with a recorded normalization scheme. In client mode, parsing, extraction, segmentation, classification proposals and encryption occur on the client, never in the storage or MCP server, and the server checks signed declarations, sizes, references and uploader authority but cannot verify plaintext extraction quality or label correctness. In server mode the storage role applies the same bounded profile and normalization scheme to the uploaded text and keeps no extraction cache. Additional formats need a separately qualified extractor with complexity limits and source-to-text mapping. Neither mode classifies content on the server. There is no server converter, plaintext indexing worker or embedding-provider fallback.}
\req{DOC-005}{R1-C}{T-06}{A text collection or conversation uses the same unit model through \code{text.entry}; no separate transcript authorization store exists. Manifests, text units, source references and policy metadata are tenant-bound catalog records, and unit IDs are not bearer permissions.}
\req{DOC-006}{R1-C}{T-06}{Document permission to enumerate or read units does not grant every unit, recipient envelope or key release. Actor, subject and delegation checks apply to the parent, selected units and envelope recipient. The original file and its separate key need explicit complete-content permission under \rr{CNT-003}; this is never inferred from a unit read. Raw, preview, thumbnail, version, manifest, envelope and plaintext routes preserve that boundary; in client mode no plaintext raw-download route exists, and in server mode the original-file read is the only plaintext raw route.}
\req{DOC-007}{R1-C}{T-06, T-09}{An explicit unit-ID request requires every requested disclosure to be allowed; otherwise it returns a generic denial. An authorized-subset operation returns only permitted units with an opaque continuation: their plaintext in server mode, or the encrypted units, their signed descriptors and separate authorized recipient envelopes in client mode. In client mode each envelope context names exactly one object or unit version and one recipient, and the qualified envelope plaintext contains exactly that unit's keyset; batching never creates an aggregate opaque multi-unit key bundle. The server checks the outer identity and grants; the recipient verifies the single-unit keyset format after unwrapping, and the server cannot inspect HPKE plaintext or prove an authorized malicious sender has not hidden other material inside it. In server mode each unit is decrypted under its own released data key. Responses expose no hidden IDs, positions, counts or placeholders. The full manifest is returned only when all contained metadata may be disclosed, and in client mode sensitive titles and text stay encrypted for the recipient.}
\req{DOC-008}{R1-C}{T-06}{Streaming-AEAD segments and HTTP byte ranges are not independent authorization units. If different spans need different rights they are split into separate units before publication.}

\subsection{Labels}\label{sec:labels}
Authorization needs a trusted classification for every unit. Authorized labeling and collection defaults suffice for the baseline; a classifier in an authorized client is optional. The server does not classify content: in client mode it cannot read it, and in server mode it does not infer labels from it. Neither an untrusted client field nor a model-generated label is authoritative.
\reqn{LAB-001}{R1-C}{T-09}{Every published unit carries a protected classification label from the tenant's label order. Publication of an unlabeled unit is rejected.}
\reqn{LAB-002}{R1-C}{T-09, T-21}{A label comes from, in order of precedence: (1) explicit labeling by an authorized principal within their delegation ceiling and at or above the collection's floor; (2) the collection or bucket default; (3) a signed proposal from a registered client-side classifier, which is advisory until accepted under (1). Client-supplied and model-generated labels are never applied directly.}
\reqn{LAB-003}{R1-C}{T-09}{A unit whose label cannot be determined takes its collection's most restrictive default, and every collection declares one.}
\reqn{LAB-004}{R1-C}{T-09}{Lowering the stored label of a unit or derived output is an explicit \code{relabel} action with its own audit record and, by default, an approval. The effective label of a derived output cannot go below its lineage maximum unless a \code{relabel} allow for that output says so (\rr{DRV-007}). Raising a label is an ordinary \code{label} action.}
\reqn{LAB-005}{R1-C}{T-09}{A label change advances the affected access revision and invalidates dependent derived-output authority under \rr{DRV-009}. Decisions, disclosure plans, server-mode key releases and stream leases bind that freshness state (\rr{DEC-006}, \rr{CNT-008}). Payload re-encryption is not required.}

\subsection{Share grants}\label{sec:grants}
A share grant lets a principal name another principal who may read a resource, within what policy allows. It is an authority fact that rules can test, never a permission by itself, and it works in both encryption modes.
\reqn{SHR-001}{R1-B}{T-24}{Recording and removing a share grant are registered, audited operations that need \code{share} on the resource. A grant names one bucket, object, document version, unit, entry or derived output; one grantee principal of the same tenant; actions from the target type's read set (\code{read}, \code{stat} and \code{discover} on an object; \code{read-units} and \code{discover} on a document version; \code{read} on a unit or entry; \code{deliver} on a derived output; \code{list} on a bucket); and an expiry no later than the grantor's own delegation. A grant to an unknown, disabled or other-tenant principal, for any other action, or beyond the grantor's delegation is rejected.}
\reqn{SHR-002}{R1-B}{T-01, T-24}{The unexpired grantees of a resource, with their granted actions, are visible to policy as \code{resource.attributes.shared_with}. Only a business rule that tests that attribute turns a grant into an allow; without such a rule a grant has no effect (default deny), and platform ceilings, labels and clearances still apply. Recording, expiring or removing a grant advances the resource's access revision, so decisions, plans, key releases and stream leases bound to the earlier state are invalidated (\rr{DEC-006}, \rr{CNT-008}).}
\reqn{SHR-003}{R1-B}{T-24}{In server mode an honored grant is enough for the storage role to decrypt for the grantee after the ordinary checks. In client mode it makes the grantee eligible under the target-recipient checks of \rr{KEY-012} but supplies no key: a client holding the data key must still publish the grantee's recipient envelope. A grant on a document reaches its units only where a rule says so, and unit checks remain per unit (\rr{DOC-006}).}

\subsection{Disclosure}
\req{DSC-001}{R1-C}{T-09}{Disclosure resolves the exact document/manifest, selected units and current authority, including transitive lineage dependencies under \rr{DRV-009}; obtains the complete parent/unit checks and, in client mode, the recipient-envelope checks; then admits that immutable plan under \rr{DAT-007}. In client mode only its authorized ciphertext, descriptors and existing envelopes are returned, and verification and decryption remain client-side. In server mode only the authorized units' plaintext is returned, each decrypted under its own released data key.}
\req{DSC-002}{R1-C}{T-09}{The server authorizes plaintext, ciphertext and envelope delivery before any of them enters a response, shared cache or downstream route. It returns server-mode plaintext only to the authenticated caller of an authorized read and never sends content to a model, reranker or embedding API on its own initiative; it never returns client-mode plaintext. A client that receives plaintext or decrypts an authorized package owns subsequent plaintext handling. Post-generation redaction and model instructions cannot substitute for the server's delivery checks or make client plaintext recallable.}
\req{DSC-003}{R1-C}{T-09}{The platform records plaintext, ciphertext and envelope delivery and server-mode key releases with exact unit/version/recipient identities and hashes, not plaintext prompts, raw keys or documents. A managed client may submit signed outcome evidence for its own optional model calls, but the server distinguishes that report from a platform-observed event and does not claim to have observed the client's plaintext use.}
\req{DSC-004}{R1-C}{T-09}{The same final unit and recipient checks apply to every disclosure surface, including preview objects, citations, notifications and subscriptions. The server emits only authorized metadata, encrypted payloads or the server-mode plaintext of authorized units. Client-generated previews and snippets are separate resources or local views; an unsupported surface is disabled.}
\req{DSC-005}{R1-C}{T-09}{Query caches are bound to tenant, authority, version and epoch and \mustnot\ reuse a result computed for a broader user for a narrower one. No raw totals or structural positions reveal denied units.}
\req{DSC-006}{R2}{T-09}{Once an authorized client receives server-mode plaintext, or a usable envelope while holding its private key, the server cannot constrain where that client sends plaintext. A destination restriction may limit which registered recipient receives content or an envelope; it does not enforce subsequent use by an arbitrary client. Any stronger managed-client promise requires a separately qualified runtime boundary and excludes that runtime's compromise. The baseline has no server-side plaintext model-proxy route.}

\subsection{Search}
\req{SRH-001}{R1-C}{T-09}{Server search is restricted to authorized catalog metadata whose visibility is explicit: tenant, permitted aliases, approved policy attributes and opaque object/version references. It performs no plaintext full-text, semantic or embedding search in either mode. After an authorized read (and, in client mode, decryption), the recipient client may search locally; its text query need not leave the client.}
\req{SRH-002}{R1-C}{T-09}{Plaintext and embedding indexes are client-side state and remain outside the server storage boundary. A qualified client protects any persistent local index with its local key custody and excludes unauthorized or undecryptable units. The server stores no plaintext snippets, embeddings or extraction cache; encrypted client indexes, if uploaded, are ordinary encrypted objects and are not server-searchable plaintext.}
\req{SRH-003}{R1-C}{T-09}{The baseline server has no remote text-search, reranking or embedding integration. A client that sends decrypted content to such a service makes an additional plaintext disclosure outside baseline storage enforcement; a managed-client extension must declare and qualify that boundary. Searchable encryption, private search and server decryption are not implied by metadata discovery.}
\reqn{SRH-004}{R1-C}{T-09}{Result values, ordering, scores and totals returned to a caller depend only on candidates that caller is authorized to see. Page shape (short or empty pages) and latency can reflect scan work over hidden candidates; they are bounded by the scan budget and tenant limits, but the platform does not claim that pagination is noninterfering with respect to hidden candidates.}

\subsection{Derived outputs and declassification}\label{sec:derived}
A client-created projection, extraction or model answer is uploaded as a new resource, encrypted under its bucket's mode. R1-C enforces labels and declared lineage for registered derived resources. These controls rely on qualified cooperating clients accurately describing their computation; the server cannot prove provenance, whether it holds ciphertext (client mode) or sees the output plaintext (server mode). Possession of plaintext by an authorized client remains a trust boundary.
\reqn{DRV-001}{R1-C}{T-09, T-21}{A managed or cooperating client that creates a derived object declares lineage naming all input units the operation received. In client mode it signs that lineage and encrypts the result before upload; in server mode it declares lineage in the authenticated upload request, bound to the caller's decision, and the storage role encrypts the result. The server validates signer or uploader authority, references, bounds and cycles. Citations do not narrow declared lineage. Because the server does not observe the computation, it cannot prove that a malicious client declared every source or that an uploaded object was independently created.}
\reqn{DRV-002}{R1-C}{T-09}{For a registered derived object, delivery requires the normal object checks (and, in client mode, envelope checks) and either (a) recipient read permission on every declared lineage unit or (b) an explicit \code{declassify} decision for that exact output, version and recipient. In client mode the owner or authorized custodian client must have created the recipient envelope, and a policy allow cannot synthesize one; in server mode the storage role decrypts for the recipient after the same checks. Declassification authorizes that delivery only and does not lower the stored label.}
\reqn{DRV-003}{R1-C}{T-09}{The managed client binds the intended recipient and conservative lineage before publishing the encrypted output. The server enforces the registered bindings without inspecting wording or inferring whether plaintext leaked. This is a provenance-based delivery rule for cooperating clients, not information-flow proof against an authorized malicious client.}
\req{DRV-004}{R2}{T-09}{Optional AI submission runs in a client or explicitly trusted managed client runtime. That client obtains permitted units (plaintext in server mode, or encrypted units and envelopes it verifies and decrypts locally in client mode) and constructs its own model request. A server may authorize the proposed recipient/model metadata, but never builds a plaintext prompt or forwards content to a model. A managed profile binds destination and parameters; baseline storage makes no claim to police arbitrary client networking.}
\reqn{DRV-005}{R2}{T-09}{The baseline rejects a request to perform server-side \code{use-in-model} without client read capability; it has no plaintext execution path for it. A future managed-client profile may distinguish an end user from a separate authorized execution recipient, but must qualify that client's custody, output encryption and recipient-specific declassification before being enabled. For client-mode content it cannot be implemented by giving the storage or MCP server a decryption key.}
\reqn{DRV-007}{R1-C}{T-09, T-21}{The effective server label of a registered derived output is at least the highest effective label in its declared lineage. A collection default acts as a floor, so the effective label is the maximum of the collection default and the lineage maximum. Only an explicit \code{relabel} allow for that output can lower it, and the effective label is the one used when the output is reopened and when a context mark is updated.}
\reqn{DRV-008}{R1-C}{T-09, T-21}{Disclosure to one recipient and relabeling are separate authorizations. A \code{declassify} allow is bound to output, version and recipient and records the label reported to a qualified managed client for flow control; absent an explicit lower label, that is the lineage maximum. Lowering the stored label for everyone is a \code{relabel} action (\rr{LAB-004}).}
\reqn{DRV-009}{R1-C}{T-09, T-13, T-21}{Derived-output authority includes its complete transitive lineage. A source access, label, lifecycle or relevant declassification change invalidates dependent decisions and leases even when output bytes and revisions are unchanged. The implementation may use dependency revisions, a conservative epoch or transactional propagation, but must qualify equivalent freshness under \rr{DAT-007}, bounded traversal and cycle rejection. New delivery recomputes effective authority; ongoing delivery follows \rr{CNT-008}. Test chained derivation and a source change during an output stream.}
\req{DRV-006}{R2}{REV}{An AI provider receiving plaintext from a client is outside the storage boundary. Client/provider retention, training and onward disclosure are not controlled by the storage cryptography in either mode. Managed-client lineage records describe declared provenance, not proof against hidden inputs, malicious client behavior or sensitive inference.}

\subsection{MCP as a later content client}
The R2 MCP service is an authenticated facade over the same content API. It preserves user and delegation authority and returns the exact versions and projections a native client can obtain. For server-mode content that is plaintext, which the storage role decrypts after the same checks, so a generic MCP client can read it. For client-mode content it is the encrypted package: a qualified client SDK or local bridge verifies the descriptors and decrypts with keys held on the client side of that boundary, a generic MCP client does not automatically understand that representation, and capability negotiation either selects the qualified format or reports it unsupported. No compatibility mode returns client-mode plaintext from the server.
\begin{longtable}{@{}P{0.22\textwidth}P{0.27\textwidth}P{0.32\textwidth}P{0.08\textwidth}@{}}
\caption{Later MCP service mappings to content operations}\label{tab:mcpops}\\
\toprule MCP operation & Required action model & Returned representation & Phase\\\midrule\endfirsthead
\toprule MCP operation & Required action model & Returned representation & Phase\\\midrule\endhead
\code{documents.search} & Authorized metadata discovery only & Permitted references and metadata; no server text snippets, embeddings or plaintext query processing & R2\\
\code{documents.read} & Parent and unit read, plus all recipient-envelope checks in client mode & Server mode: plaintext of the permitted units. Client mode: selected ciphertext, signed unit descriptors and existing recipient envelopes; local decryption & R2\\
\code{text.read} & Collection/entry and exact-version read, plus recipient checks in client mode & The same unit representation as \code{documents.read}; no unguarded text endpoint & R2\\
\code{resources/list}, \code{resources/read} & Same discovery and read checks as tools & Registry-bound references and the same representation as the native read; no direct backend or file access & R2\\
\code{documents.export} & Complete mapping of Table~\ref{tab:keyops} & Client mode: ciphertext and envelopes for an approved compatible recipient, with no server rewrapping or decryption. Server mode: plaintext of a permitted complete export & R2\\
\bottomrule
\end{longtable}
Optional AI assistance is a client workflow after an authorized read. The baseline MCP service exposes no \code{ai.ask} operation that receives or constructs plaintext context; any managed-client model integration needs its separately qualified authority and disclosure boundary.

\subsection{Sizing and scaling}
\reqn{SIZ-001}{R1-B}{T-22}{Qualification reports measured object-service capacity in R1-B for both modes: server-mode encryption and decryption throughput and data-key wrap and unwrap rates, and client-mode ciphertext transfer with client encrypt/decrypt costs; then unit and recipient-envelope fan-out for the reference documents in R1-C. R2 adds the independently deployed MCP service. Measurements separate server payload cryptography, ciphertext throughput and signature checks from client HPKE wrapping/unwrapping and payload AEAD work.}
\req{SIZ-002}{R1-B}{T-22}{The \code{storage} role scales independently with bounded requests, payload streams, server-mode encryption work, upload/read concurrency, buffer memory, signature-verification work and tenant byte budgets; the \code{keys} role scales with key releases. Decision replicas and backend capacity are qualified with \rr{AVL-007}; adding frontends does not imply unlimited catalog throughput. R1-C adds bounded unit and envelope fan-out and metadata discovery. The R2 \code{mcp} role has its own replica set and identity and calls the same content API. Client compute and local indexes scale at the client; no server permission database, evaluator or plaintext cache is added for MCP.}
\begin{longtable}{@{}P{0.28\textwidth}P{0.44\textwidth}P{0.18\textwidth}@{}}
\caption{Sizing quantities (formulas; values are measured, not assumed)}\label{tab:sizing}\\
\toprule Quantity & Measurement or formula, with $u$ units per document and $c$ search candidates & Release\\\midrule\endfirsthead
\toprule Quantity & Measurement or formula, with $u$ units per document and $c$ search candidates & Release\\\midrule\endhead
Object-service capacity & Object-size distribution, concurrent streams, read/write throughput, catalog commit rate, backend latency and bounded memory per request & R1-B\\
Decision batches to read a whole document & $\lceil k/50 \rceil$, where $k$ includes the parent, $u$ unit checks and every required recipient-envelope check & R1-C\\
Decision batches per metadata page & $\lceil c/50 \rceil$ before any prefilter (\rr{CNT-007}) & R1-C\\
Blobs, catalog rows per ingest & $u$ blobs, $u$ unit rows, 1 manifest & R1-C\\
Recipient envelopes (client mode) & Up to $u \times r$ client-created HPKE envelopes for $r$ authorized recipients; server stores and authorizes delivery only & R1-C\\
Client cryptographic work (client mode) & $u$ independent payload encryptions at ingest; selected-unit envelope unwraps and payload decryptions at recipients & R1-C\\
Server cryptographic work (server mode) & $u$ data-key wraps and payload encryptions at ingest; one key release and payload decryption per delivered unit & R1-C\\
Durable receipts per publication & Backend-defined; 2 per immutable blob with the custom store & R1-B, R3\\
Reference documents & $u = 10$, $u = 1{,}000$ and $u = 10{,}000$ & R1-C\\
\bottomrule
\end{longtable}
